% Source-aligned English for printed p.214.
\clearpage
\SourcePage{214}

\LayerTarget{tohoku:definition:equivariant-o-homomorphisms}
\noindent the space on which $A$ and $B$ are considered (which also gives
meaning to the symbol $\operatorname{Hom}_{O,G}(U;A,B)$ if $U$ is a subset of
$X$ stable under $G$). Since the sheaf $\mathbf{Hom}_O(A,B)$ is a
$G$-sheaf, the group
$\operatorname{Hom}_O(A,B)=\Gamma_X\mathbf{Hom}_O(A,B)$ is a $G$-module, and
we have the formulas

\LayerTarget{tohoku:equation:5.6.1}
\[
\tag{5.6.1}
\operatorname{Hom}_{O,G}(A,B)
=\Gamma_X^G\mathbf{Hom}_O(A,B)
=\Gamma^G\operatorname{Hom}_O(A,B).
\]

\LayerTarget{tohoku:definition:equivariant-hom-sheaf}
Put:

\LayerTarget{tohoku:equation:5.6.2}
\[
\tag{5.6.2}
\mathbf{Hom}_{O,G}(A,B)=(\mathbf{Hom}_O(A,B))^G
\]

\noindent thus $\mathbf{Hom}_{O,G}(A,B)$ is the sheaf on $Y=X/G$ whose
sections over the open set $U\subset Y$ form the group
$\operatorname{Hom}_{O,G}(f^{-1}U;A,B)$ of $G$-$O$-homomorphisms from
$A|f^{-1}(U)$ to $B|f^{-1}(U)$. Put:

\LayerTarget{tohoku:definition:four-equivariant-hom-functors}
\LayerTarget{tohoku:equation:5.6.3}
\[
\tag{5.6.3}
\begin{aligned}
\mathbf h_{O,A}(B)&=\mathbf{Hom}_O(A,B),&
h_{O,A}(B)&=\operatorname{Hom}_O(A,B),\\
\mathbf h_{O,G,A}(B)&=\mathbf{Hom}_{O,G}(A,B),&
h_{O,G,A}(B)&=\operatorname{Hom}_{O,G}(A,B),
\end{aligned}
\]

\LayerTarget{tohoku:relation:equivariant-hom-functors}
\noindent thus defining, for fixed $A$, four left exact functors, the latter
two being related to the former two (already encountered in Chapter IV) by
the functorial formulas

\LayerTarget{tohoku:equation:5.6.4}
\[
\tag{5.6.4}
\mathbf h_{O,G,A}=f_*^G\mathbf h_{O,A}
\]

\LayerTarget{tohoku:equation:5.6.5}
\[
\tag{5.6.5}
h_{O,G,A}=\Gamma_Y\mathbf h_{O,G,A}
=\Gamma^G h_{O,A}=\Gamma_X^G\mathbf h_{O,A}.
\]

\LayerTarget{tohoku:definition:equivariant-ext}
We shall put:

\LayerTarget{tohoku:equation:5.6.6}
\[
\tag{5.6.6}
\operatorname{Ext}_{O,G}^n(X;A,B)=R^n h_{O,G,A}(B),
\]

\LayerTarget{tohoku:equation:5.6.7}
\[
\tag{5.6.7}
\mathbf{Ext}_{O,G}^n(A,B)=R^n\mathbf h_{O,G,A}(B).
\]

\LayerTarget{tohoku:property:equivariant-ext-functoriality}
Thus, for two $G$-$O$-modules $A,B$, the
$\operatorname{Ext}_{O,G}^n(X;A,B)$ are abelian groups forming a universal
cohomological functor in $B$ and reducing to
$\operatorname{Hom}_{O,G}(X;A,B)$ in dimension 0, whereas the
$\mathbf{Ext}_{O,G}^n(A,B)$ are abelian sheaves on $Y$, forming a universal
cohomological functor in $B$ and reducing to $\mathbf{Hom}_{O,G}(A,B)$ in
dimension 0. Moreover, the $\operatorname{Ext}_{O,G}^n(X;A,B)$ are precisely
the general $\operatorname{Ext}$ groups in the category $C^{O(G)}$. (Our
definitions are, of course, justified by Proposition 5.1.1.) Since
$\operatorname{Hom}_{O,G}(A,B)$ and hence also
$\mathbf{Hom}_{O,G}(A,B)$ is a contravariant functor exact in $A$ whenever
$B$ is an injective object of $C^{O(G)}$, it follows (cf. 2.3) that the
$\operatorname{Ext}_{O,G}^n(X;A,B)$, respectively the
$\mathbf{Ext}_{O,G}^n(A,B)$, also form a cohomological functor contravariant
in $A$ (for fixed $B$). Finally, one proves easily, as in 4.2 and 5.2 (which
are special cases), that

\LayerTarget{tohoku:equation:5.6.8}
\[
\tag{5.6.8}
\mathbf{Ext}_{O,G}^n(A,B)|U
=\mathbf{Ext}_{O,G}^n(A|f^{-1}(U),B|f^{-1}(U))
\]

\noindent (\textit{local nature of} $\mathbf{Ext}$ \textit{with respect to
the space} $Y$). It follows that there is a statement analogous to the
corollary to Proposition 5.2.1, which, when $G$ is a \textit{discontinuous}
group of homeomorphisms (cf. 5.3), permits reduction to the case in which $G$
is finite. Repeating the arguments of 4.1, one then proves that the natural
homomorphisms

\LayerTarget{tohoku:comparison:equivariant-hom-stalk-to-fiber}
\LayerTarget{tohoku:equation:5.6.9}
\[
\tag{5.6.9}
\mathbf{Hom}_{O,G}(A,B)(y)
\longrightarrow \operatorname{Hom}_{U_x}(A(x),B(x))
\qquad (f(x)=y)
\]

\noindent (where $U_x$ is the ring generated by $O(x)$ and $Z(G_x)$, considered
in 5.1) are

\TohokuContinuation{215}
